Recent evidence associates generative-AI exposure with weaker employment among workers in their early twenties. I reconstruct the public monthly CPS comparison and ask which occupational contrasts support it. The estimand is the post-December-2022 change in the log conditional-mean employment-stock ratio of ages 22--25 to 26--65 in the highest versus lowest GPT-4-exposure quintile, conditional on intermediate quintiles and preperiod software exposure. Although nominal support contains 468 occupations, only four broad families span both tails; their 29 occupations represent 5.03 percent of preperiod stock, and direct-tail information is equivalent to 6.2 occupations. The pooled estimate is $-0.132$ log points. Broad occupation-family paths reduce it to $-0.022$, but the direct-tail interval cannot resolve a within-family gradient. An annual ACS extension narrows person-sampling uncertainty without creating the missing occupational comparisons. Computer-use and family conditioning move the coefficient in opposite directions, preperiod dynamics reject parallel evolution, and linked-CPS estimates do not isolate a flow mechanism. The public-CPS association is supported predominantly by cross-family comparisons; it is not an identified causal AI effect or a replication of employer hiring results.
